\usepackage[round,authoryear]{natbib}
\usepackage{paralist}
\usepackage{pgfpages}
\usepackage{mathtools}
\usepackage{amsmath}
\usepackage{amssymb}
\usepackage{xspace}        
\usepackage{graphicx}
\usepackage{listings}

\bibliographystyle{jss}
%\bibliographystyle{apalike}
%\bibliographystyle{dcu}

\usefonttheme[onlymath]{serif}
% use article-like math letters, from http://tex.stackexchange.com/questions/34265/how-to-get-beamer-math-to-look-like-article-math

%%% notation for ARMA models
\newcommand\ar{\phi}
\newcommand\ma{\psi}
\newcommand\AR{\Phi}
\newcommand\MA{\Psi}

%%% notation for linear Gaussian POMP models
\newcommand\matA{\mathbb{A}}
\newcommand\matB{\mathbb{B}}
\newcommand\matH{\mathbb{H}}
\newcommand\covmatX{\mathbb{U}}
\newcommand\covmatY{\mathbb{V}}

\newcommand\prob[1]{\mathbb{P}\left[{#1}\right]}
\newcommand\expect[1]{\mathbb{E}\left[{#1}\right]}
\newcommand\var[1]{\mathrm{Var}\left[{#1}\right]}
\newcommand\myvar{\mathrm{Var}}
\newcommand\dist[2]{\mathrm{#1}\left(#2\right)}
\newcommand\lik{\mathcal{L}}
\newcommand\loglik{\ell}
\newcommand\R{\mathbb{R}}
\newcommand\Rzero{\mathfrak{R}_0}
\newcommand\argmax{\mathop{\mathrm{argmax}}}
\newcommand\argmin{\mathop{\mathrm{argmin}}}

\newcommand\code[1]{\texttt{#1}}
\newcommand\package[1]{\textbf{#1}}
\newcommand\data[1]{#1^*}
%\newcommand\data[1]{#1}
%\newcommand\estimate[1]{\data{#1}}
\newcommand\estimate[1]{#1}
\newcommand\predict[1]{#1^{P}}
\newcommand\residual{r}

\newcommand\link[2]{\href{#1}{#2}}
\newcommand{\doi}[1]{\link{https://doi.org/#1}{\texttt{doi:~{#1}}}}

\newcommand{\scinot}[2]{#1{\times}10^{#2}}
\newcommand{\dd}[1]{\mathrm{d}{#1}}
\newcommand{\pd}[3][]{%
  \def\ord{#1} \ifx\ord\empty%
  \frac{\partial{#2}}{\partial{#3}}%
  \else \frac{\partial^{#1}{#2}}{\partial{#3}^{#1}}%
  \fi
}
\newcommand{\deriv}[3][]{%
  \def\ord{#1} \ifx\ord\empty%
  \frac{\dd{#2}}{\dd{#3}}%
  \else \frac{\dd^{#1}{#2}}{\dd{#3}^{#1}}%
  \fi
}

\newcommand\Rlanguage{\textsf{R}\xspace}

\newcommand\question{{\bf Question}}
\newcommand\mysolution{{\bf Solution}}
\newcounter{Qcounter}
\newcommand\myquestion{{\stepcounter{Qcounter} \bf Question \CHAPTER.\theQcounter}}
\newcounter{Ecounter}
\newcommand\myexercise{{\stepcounter{Ecounter} \bf Exercise \CHAPTER.\theEcounter}}

\newcommand\myexample{{\bf Example}}
\newcommand\mydot{{\,\cdot\,}}
\newcommand\myref[1]{\m{#1}}

\newcommand\Rspace{\mathcal{R}}

\renewcommand\vec[1]{\boldsymbol{\mathrm{#1}}}
\newcommand\vect[1]{\vec{#1}}
\newcommand\mat[1]{\mathbb{#1}}
\newcommand\pr{\mathbb{P}}
\newcommand\E{\mathbb{E}}

\newcommand\profileloglik[1]{\ell^\mathrm{profile}_#1}
\newcommand\Real{\mathbb{R}}

\newcommand\bi{\begin{itemize}}
\newcommand\ei{\end{itemize}}

\newcommand\normal{\mathrm{normal}}

\newcommand\iid{iid}
\newcommand\MVN{\mathrm{MVN}}
\newcommand\SE{\mathrm{SE}}

\newcommand\pval{\mathrm{pval}}
% \newcommand\var{\mathrm{Var}}
\newcommand\sd{\mathrm{SD}}
\newcommand\sdSample{\mathrm{sd}}
\newcommand\varSample{\mathrm{var}}
\newcommand\cov{\mathrm{Cov}}
\newcommand\covSample{\mathrm{cov}}
\newcommand\corSample{\mathrm{cor}}
\newcommand\cor{\mathrm{Cor}}
\newcommand\given{{\, | \,}}
\newcommand\param{\,;}
\newcommand\params{\param}
\newcommand\equals{{ \, = \, }}
\newcommand\transpose{{\raisebox{0.5mm}{\mbox{\scriptsize \textsc{t}}}}}
%\newcommand\transpose{\scriptsize{T}}
\newcommand\mycolon{{\hspace{0.5mm}:\hspace{0.5mm}}}
%\newcommand\mycolon{\,{:}\,}
\newcommand\myemph[1]{{\textbf{#1}}}
\newcommand\mymathenv[1]{\textcolor{blue}{#1}}


%% following a style where all math is in blue
%% not currently used 
\newcommand\mymath[1]{\begin{math}\textcolor{blue}{#1}\end{math}}
\newcommand\m[1]{\mymath{#1}}
\newcommand\mydisplaymath[1]{\begin{displaymath}\textcolor{blue}{#1}\end{displaymath}}
\newcommand\myeqnarray[1]{\textcolor{blue}{\begin{eqnarray*}#1 \end{eqnarray*}}}
\newcommand\myspace{\quad}
\newcommand\altdisplaymath[1]{\vspace{1mm}\textcolor{blue}{\begin{math}\displaystyle #1 \end{math}}\vspace{1mm}}

%%%%%%%%%%%%%%%%%


\setlength{\parskip}{0mm}
\setlength{\parindent}{0mm}
\lstset{language=C}

\newcommand\negBeforeCode{}
\newcommand\negAfterCode{}

\mode<beamer>{\usetheme{AnnArbor}}
\mode<beamer>{\setbeamertemplate{footline}}
\mode<beamer>{\setbeamertemplate{footline}[frame number]}
\mode<beamer>{\setbeamertemplate{frametitle continuation}[from second][\insertcontinuationcountroman]}
\mode<beamer>{\setbeamertemplate{navigation symbols}{}}

\mode<handout>{\pgfpagesuselayout{2 on 1}[letterpaper,border shrink=5mm]}


\newcommand\eqspace{\hspace{3mm}}
\newcommand\eqvspace{\vspace{1mm}}
\newcommand\negListSpace{\hspace{-4mm}}

\newcommand\ev{u}

\mode<beamer>{\usetheme{AnnArbor}}
\mode<beamer>{\setbeamertemplate{footline}}
\mode<beamer>{\setbeamertemplate{footline}[frame number]}
\mode<beamer>{\setbeamertemplate{frametitle continuation}[from second][\insertcontinuationcountroman]}
\mode<beamer>{\setbeamertemplate{navigation symbols}{}}

\mode<handout>{\pgfpagesuselayout{2 on 1}[letterpaper,border shrink=5mm]}

\newcommand\CHAPTER{7}
% \newcommand\answer[2]{\textcolor{blue}{#2}} % to show answers
% \newcommand\answer[2]{\textcolor{red}{#2}} % to show answers
 \newcommand\answer[2]{#1} % to show blank space
\usepackage{bbm} % for blackboard bold 1


\title{\vspace{2mm} \link{https://ionides.github.io/531w24/}{Modeling and Analysis of Time Series Data}\\ \vspace{2mm}
Chapter \CHAPTER: Introduction to time series analysis in the frequency domain}
\author{Edward L. Ionides}
\date{}

\setbeamertemplate{footline}[frame number]




\begin{document}

\maketitle

\mode<article>{\tableofcontents}

\mode<presentation>{
  \begin{frame}{Outline}
    \tableofcontents
  \end{frame}
}




\section{Frequency components}
%\subsection{}

\begin{frame}{Frequency components of a time series}

\begin{enumerate}
\item A time series dataset (like any other sequence of numbers) can be written as a sum of sine and cosine functions with varying frequencies.
\item This is called the \myemph{Fourier representation} or \myemph{Fourier transform} of the data.
\item The coefficients corresponding to the sine and cosine at each frequency are called \myemph{frequency components} of the data.
\item   Looking at which frequencies have large and small components can help to identify appropriate models.
\item Looking at the frequency components present in our models can help to assess whether they are doing a good job of describing our data.

\end{enumerate}

\end{frame}

\subsection{Eigenvalues and eigenvectors of a stationary covariance matrix}

\begin{frame}[fragile]{What is the spectrum of a time series model?}

\mode<presentation>{
\vspace{-1mm}
}

\bi
\item  We begin by reviewing eigenvectors and eigenvalues of covariance matrices. This eigen decomposition also arises elsewhere in statistics, e.g. principle component analysis.

\item  A univariate time series model is a vector-valued random variable $Y_{1:N}$ which we suppose has a covariance matrix $V$ which is an $N\times N$ matrix with entries $V_{mn}=\cov(Y_m,Y_n)$.

\item  $V$ is a non-negative definite symmetric matrix, and therefore has $N$ non-negative eigenvalues $\lambda_1,\dots,\lambda_N$ with corresponding eigenvectors $\ev_1,\dots,\ev_N$ such that
\begin{equation} V \ev_n = \lambda_n \ev_n.\end{equation}

\item  A basic property of these eigenvectors is that they are orthogonal, i.e.,
\begin{equation} \ev_m^\transpose \ev_n = 0 \mbox{ if $m\neq n$}.\end{equation}

\item  We may work with \myemph{normalized} eigenvectors that are scaled such that $\ev_n^\transpose \ev_n = 1$.

\ei

\end{frame}

\begin{frame}[fragile]

\bi
\item  We can also check that the components of $Y$ in the directions of different eigenvectors are uncorrelated.

\item Since $\cov(AY,BY)=A\cov(Y,Y)B^\transpose$, we have
\begin{eqnarray*}
\cov(\ev_m^\transpose Y, \ev_n^\transpose Y) &=& \ev_m^\transpose \cov(Y,Y) \ev_n
\\
&=& \ev_m^\transpose V \ev_n
\\
&=&\lambda_n \ev_m^\transpose \ev_n
\\
&=& \left\{
  \begin{array}{cc} 
    \lambda_n & \mbox{if $m=n$} \\
    0 & \mbox{if $m\neq n$}
  \end{array}
\right.
\end{eqnarray*}
For the last equality, we have supposed that the eigenvectors are normalized.

\item  If we knew $V$, we could convert the model to a representation where the observable random variables are uncorrelated. 

\item  Transforming the data into its components in the directions of the eigenvectors of the model allows us to use an uncorrelated model. In the Gaussian case, we have independence.

\ei

\end{frame}

\begin{frame}[fragile]

\frametitle{Eigenvectors for the covariance matrix of an AR(1) model with $N=100$ and $\phi=0.8$}

\vspace{-1mm}

\begin{knitrout}\small
\definecolor{shadecolor}{rgb}{0.969, 0.969, 0.969}\color{fgcolor}\begin{kframe}
\begin{alltt}
\hldef{N} \hlkwb{<-} \hlnum{100}\hldef{;  phi} \hlkwb{<-} \hlnum{0.8}\hldef{;  sigma} \hlkwb{<-} \hlnum{1}
\hldef{V} \hlkwb{<-} \hlkwd{matrix}\hldef{(}\hlnum{NA}\hldef{,N,N)}
\hlkwa{for}\hldef{(m} \hlkwa{in} \hlnum{1}\hlopt{:}\hldef{N)} \hlkwa{for}\hldef{(n} \hlkwa{in} \hlnum{1}\hlopt{:}\hldef{N) V[m,n]}\hlkwb{<-}\hldef{sigma}\hlopt{^}\hlnum{2}\hlopt{*}\hldef{phi}\hlopt{^}\hlkwd{abs}\hldef{(m}\hlopt{-}\hldef{n)}\hlopt{/}\hldef{(}\hlnum{1}\hlopt{-}\hldef{phi}\hlopt{^}\hlnum{2}\hldef{)}
\hldef{V_eigen} \hlkwb{<-} \hlkwd{eigen}\hldef{(V,}\hlkwc{symmetric}\hldef{=}\hlnum{TRUE}\hldef{)}
\hlkwd{matplot}\hldef{(V_eigen}\hlopt{$}\hldef{vectors[,}\hlnum{1}\hlopt{:}\hlnum{5}\hldef{],}\hlkwc{type}\hldef{=}\hlsng{"l"}\hldef{)}
\hlkwd{matplot}\hldef{(V_eigen}\hlopt{$}\hldef{vectors[,}\hlnum{6}\hlopt{:}\hlnum{9}\hldef{],}\hlkwc{type}\hldef{=}\hlsng{"l"}\hldef{)}
\end{alltt}
\end{kframe}
\end{knitrout}

\vspace{-2mm}

\begin{knitrout}\small
\definecolor{shadecolor}{rgb}{0.969, 0.969, 0.969}\color{fgcolor}

{\centering \includegraphics[width=11cm]{tmp/figure/eigen-1} 

}


\end{knitrout}

\end{frame}

\begin{frame}[fragile]

\frametitle{Eigenvalues for the covariance matrix of an AR(1) model with $N=100$ and $\phi=0.8$}

\bi
\item  We see that the eigenvectors, plotted as functions of time, look like sine wave oscillations.

\item  The eigenvalues are
\begin{knitrout}\small
\definecolor{shadecolor}{rgb}{0.969, 0.969, 0.969}\color{fgcolor}\begin{kframe}
\begin{alltt}
\hlkwd{round}\hldef{(V_eigen}\hlopt{$}\hldef{values[}\hlnum{1}\hlopt{:}\hlnum{9}\hldef{],}\hlnum{2}\hldef{)}
\end{alltt}
\begin{verbatim}
[1] 24.59 23.44 21.73 19.70 17.57 15.51 13.61 11.91 10.42
\end{verbatim}
\end{kframe}
\end{knitrout}

\item  We see that the eigenvalues are decreasing. For this model, the components of $Y_{1:N}$ with highest variance correspond to long-period oscillations.

\item  Are the sinusoidal eigenvectors a special feature of this particular time series model, or something more general?

\ei

\end{frame}


\begin{frame}[fragile]
\frametitle{The eigenvectors for a long stationary time series model}

\mode<presentation>{
\vspace{-2.5mm}
}

\bi
\item  Suppose $\{Y_n,-\infty<n<\infty\}$ has a stationary autocovariance function $\gamma_h$. Write $\Gamma$ for the infinite array with entries
\begin{equation} \Gamma_{m,n} = \gamma_{m-n} \quad \mbox{for all integers $m$ and $n$}.\end{equation}
\item To focus on concepts over technical details, we assume infinite sums converge and order of summation can be exchanged, so infinite arrays behave like finite matrices.
\item  An eigenvector for $\Gamma$ is a sequence $\ev=\{\ev_n, -\infty<n<\infty\}$ with corresponding eigenvalue $\lambda$ such that
\begin{equation}\Gamma \ev = \lambda \ev,\end{equation}
or, writing out the matrix multiplication explicitly,

\mode<presentation>{
\vspace{-3mm}
}

\begin{equation}\sum_{n=-\infty}^\infty \Gamma_{m,n} \ev_n = \lambda \ev_m\quad \mbox{for all $m$}.\end{equation}

\mode<presentation>{
\vspace{-1mm}
}

\item  We look for a sinusoidal solution, $\ev_n = e^{2\pi i\omega n}$, where $\omega$ is cycles per unit time.

\ei

\end{frame}

\begin{frame}[fragile]

\vspace{-3mm}

\begin{eqnarray*} \textstyle
\sum_{n=-\infty}^\infty \Gamma_{m,n} \ev_n 
&=&
\textstyle
\sum_{n=-\infty}^\infty \gamma_{m-n} \ev_n 
\\
&=&
\textstyle
\sum_{h=-\infty}^\infty \gamma_{h}  \ev_{m-h} \quad \mbox{setting $h=m-n$}
\\
&=&
\textstyle
\sum_{h=-\infty}^\infty \gamma_{h}  e^{2\pi i\omega(m-h)}
\\
&=&
\textstyle
e^{2\pi i\omega m} \sum_{h=-\infty}^\infty \gamma_{h}  e^{-2\pi i\omega h}
\\
&=&
\textstyle
\ev_m \lambda(\omega) \hspace{3mm} \mbox{ for } \lambda(\omega)= \sum_{h=-\infty}^\infty \gamma_{h}  e^{-2\pi i\omega h}
\end{eqnarray*}

\myquestion. Why does this calculation show that 
$\ev_n(\omega) = e^{2\pi i\omega n}$
is an eigenvector for $\Gamma$ for any choice of $\omega$.

\answer{\vspace{40mm}}{todo}

\end{frame}
\begin{frame}[fragile]

\bi
\item 
The eigenvalue at frequency $\omega$ is
\begin{equation}
\label{eq:ft1}
\lambda(\omega)= \sum_{h=-\infty}^\infty \gamma_{h} \,  e^{-2\pi i\omega h}\end{equation}
\item Viewed as a function of $\omega$, this is called the \myemph{spectral density function}.

\item $\lambda(\omega)$ is the \myemph{Fourier transform} of $\gamma_h$.

\item An integral version of (\ref{eq:ft1}) is used in applied math and engineering:
\begin{equation}
\label{eq:ft2}
\lambda(\omega) = \int_{-\infty}^{\infty} \gamma(x) \, e^{-2\pi i\omega x}\, dx.
\end{equation}
\item We obtain (\ref{eq:ft1}) from (\ref{eq:ft2}) when $\gamma(h)$ has a point mass $\gamma_h$ when $h$ is an integer, and $\gamma(x)=0$ for non-integer $x$.
\ei

\end{frame}

\begin{frame}

\bi
\item  It was convenient to do this calculation with complex exponentials. However, writing
\begin{equation} e^{2\pi i\omega n} = \cos(2\pi\omega n) + i \sin(2\pi\omega n),\end{equation}
and noting that $\gamma_h$ is real, we see that the real and imaginary parts of $\lambda(\omega)= \sum_{h=-\infty}^\infty \gamma_{h}  e^{-2\pi i\omega h}$ give us two real eigenvectors, $\cos(2\pi\omega n)$ and $\sin(2\pi\omega n)$.
\ei

\vspace{2mm}

\myquestion. Review: how would you demonstrate the correctness of the identity
$e^{2\pi i\omega} = \cos(2\pi\omega)+i\,\sin(2\pi\omega)$.

\answer{\vspace{20mm}}{todo}

\end{frame}

\begin{frame}[fragile]

\bi

\item  Assuming that this computation for an infinite sum represents a limit of increasing dimension for finite matrices, we have found that the eigenvectors for any long, stationary time series model are approximately sinusoidal.

\item  For the finite time series situation, we only expect $N$ eigenvectors for a time series of length $N$. We have one eigenvector for $\omega=0$, two eigenvectors corresponding to sine and cosine functions with frequency
\begin{equation}\omega_{n} = n/N, \mbox{ for $0<n<N/2$},\end{equation}
and, if $N$ is even,  a final eigenvector with frequency
\begin{equation}
\omega_{(N/2)} = 1/2.
\end{equation}

\item  These sine and cosine vectors are the \myemph{Fourier basis}.

\item The time series  $\data{y}_{1:N}$ is the \myemph{time domain} representation of the data. Transforming to the Fourier basis gives the \myemph{frequency domain} representation.

\ei

\end{frame}

\section{The Fourier transform}

\begin{frame}[fragile]
\frametitle{Frequency components and the Fourier transform}

\bi
\item  The \myemph{frequency components} of $Y_{1:N}$ are the components in the directions of these eigenvectors, given by
\begin{eqnarray*}
C_n &=& \frac{1}{\sqrt{N}}\sum_{k=1}^N Y_k\cos(2\pi \omega_n k) \mbox{ for $0\le n\le N/2$},
\\
S_n &=& \frac{1}{\sqrt{N}}\sum_{k=1}^N Y_k\sin(2\pi \omega_n k) \mbox{ for $1\le n\le N/2$}.
\end{eqnarray*}

\item  Similarly, the \myemph{frequency components} of data $\data{y_{1:N}}$ are 
\begin{eqnarray*}
{c_n} &=& \frac{1}{\sqrt{N}}\sum_{k=1}^N \data{y_k}\cos(2\pi \omega_n k) \mbox{ for $0\le n\le N/2$},
\\
{s_n} &=& \frac{1}{\sqrt{N}}\sum_{k=1}^N \data{y_k}\sin(2\pi \omega_n k) \mbox{ for $1\le n\le N/2$}.
\end{eqnarray*}

\ei

\end{frame}

\begin{frame}[fragile]
\bi
\item  The frequency components of the data can be written as real and imaginary parts of the \myemph{discrete Fourier transform},
\begin{eqnarray*}
{d_n} &=& \frac{1}{\sqrt{N}} \sum_{k=1}^N \data{y_k} e^{-2\pi i k n/N}
\\
&=&{c_n} - i{s_n}
\end{eqnarray*}

\item  The normalizing constant of $1/\sqrt{N}$ is convenient for a central limit theorem.

\item Various choices about signs and factors of $2\pi$, $\sqrt{2\pi}$ and $\sqrt{N}$ can be made in the definition of the Fourier transform. For example, the \code{fft} command in {\Rlanguage} does not include this constant. 

\item  \code{fft} is an implementation of the fast Fourier transform algorithm, which enables computation of all the frequency components with order $N\log(N)$ computation. Computing the frequency components may appear to require a matrix multiplication involving order $N^3$ additions and multiplications. When $N=10^5$ or $N=10^6$ this difference becomes important!

\ei

\end{frame}

\begin{frame}[fragile]

\bi

\item  The first frequency component, $C_0$, is a special case, since it has mean $\mu=\E[Y_n]$ whereas the other components have mean zero.

\item  In practice, we subtract a mean before computing the frequency components, which is equivalent to removing the frequency component for frequency zero.

\item  The frequency components $(C_{0:N/2},S_{1:N/2})$ are asymptotically uncorrelated. They are constructed as a sum of a large number of terms, with the usual  $1/\sqrt{N}$ scaling for a central limit theorem. So, it may not be surprising that a central limit theorem applies, giving asymptotic justification for the following normal approximation. 

\item  Moving to the frequency domain (i.e., transforming the data to its frequency components) has \myemph{decorrelated} the data. Statistical techniques based on assumptions of independence are appropriate when applied to frequency components.
\ei

\end{frame}

\subsection{ A normal approximation}

\begin{frame}[fragile]{Normal approximation for the frequency components}

\bi

\item  $(C_{1:N/2},S_{1:N/2})$ are approximately independent, mean zero, Normal random variables with
\begin{equation} \myvar(C_n) = \myvar(S_n) \approx 1/2 \lambda(\omega_n).\end{equation}

\item  $C_0\big/ \sqrt{N}$ is approximately Normal, mean $\mu$, independent of $(C_{1:N/2},S_{1:N/2})$, with
\begin{equation}\myvar(C_0\big/ \sqrt{N}) \approx \lambda(0)\big/ N.\end{equation}



\item  It follows from the normal approximation that, for $1\le n\le N/2$,
\begin{equation}
\label{eq:cs:chi-squared}
C_n^2 + S_n^2 \approx \lambda(\omega_n) \frac{\chi^2_2}{2},\end{equation}
where $\chi^2_2$ is a chi-squared random variable on two degrees of freedom.

\item  Taking logs, we have
\begin{equation}
\label{eq:cs:log-chi-squared}
\log\big(C_n^2 + S_n^2 \big) \approx \log \lambda(\omega_n) + \log(\chi^2_2/2).\end{equation}

\ei
\end{frame}

\section{The periodogram to estimate the spectral density}

\begin{frame}{The periodogram}

\bi
\item  The chi-squared property in (\ref{eq:cs:chi-squared}) motivates the \myemph{periodogram},
\begin{equation} I_n = {{c_n}}^2 + {{s_n}}^2 = \big|  {d_n}\big|^2\end{equation}
as an estimator of the spectral density. 

\item  From (\ref{eq:cs:log-chi-squared}), $\log I_n$ is as an estimator of the log spectral density with a convenient statistical property: asymptotically independent, identically distributed errors at each Fourier frequency.

\item  Therefore, a signal-plus-white-noise model is appropriate for estimating the log spectral density using the log periodogram.

\item The periodogram is an \myemph{inconsistent estimator} of the spectrum. We can smooth the periodogram to borrow strength between nearby frequencies.

\ei

\end{frame}


\begin{frame}[fragile]

\frametitle{Interpreting the spectral density as a power spectrum}

\bi

\item  The power of a wave is proportional to the square of its amplitude. 

\item  The spectral density gives the mean square amplitude of the components at each frequency, and therefore gives the expected power.

\item  The spectral density function can therefore be called the \myemph{power spectrum}.

\ei

\end{frame}

\begin{frame}
\myquestion. Consider the AR(1) model, $\phi(B)Y_n = \epsilon_n$ with $\phi(B)=1-\phi_1B$ and $\epsilon_n \sim \mathrm{WN}(\sigma^2)$, i.e., white noise with variance $\sigma^2$. Show that the spectrum of $Y$ is
\begin{equation}
\lambda(\omega)=\frac{\sigma^2}{\big|\phi\big( e^{2\pi i \omega}\big)\big|^2}
= \frac{\sigma^2}{1+\phi_1^2 - 2\phi_1\cos(2\pi\omega)}.
\end{equation}

\answer{\vspace{60mm}}{todo}

\end{frame}

\begin{frame}{ARMA models have a rational spectrum}
\bi
\item The calculation for the AR(1) model generalizes. We give the result without proof.
\item Let $Y_n$ be an ARMA(p,q) model, $\ar(B)Y_n = \ma(B)\epsilon_n$ with $\epsilon_n\sim \mathrm{WN}(\sigma^2)$. The spectrum of $Y$ is
\begin{equation} \label{eq:arma_spec}
\lambda(\omega) = \sigma^2 \left| \frac{\psi\big( e^{2\pi i \omega} \big)}{\phi\big( e^{2\pi i \omega} \big)} \right|^2.
\end{equation}
\item The so-called \myemph{rational spectrum} of ARMA models is computationally convenient.
\item A stationary, causal ARMA model cannot have roots on the unit circle. If a root approaches the unit circle, the denominator in (\ref{eq:arma_spec}) becomes close to zero.
\item The special case of $\phi(x)=\psi(x)=1$ gives $\lambda(\omega)=\sigma^2$. \myemph{White noise has a constant spectrum}, matching the analogy that white light has uniform intensity across the visible light spectrum.

\ei

\end{frame}

\section{Frequency domain data analysis}

\begin{frame}[fragile]

\frametitle{Michigan winters revisited: Frequency domain methods}



\begin{knitrout}\small
\definecolor{shadecolor}{rgb}{0.969, 0.969, 0.969}\color{fgcolor}\begin{kframe}
\begin{alltt}
\hldef{y} \hlkwb{<-} \hlkwd{read.table}\hldef{(}\hlkwc{file}\hldef{=}\hlsng{"ann_arbor_weather.csv"}\hldef{,}\hlkwc{header}\hldef{=}\hlnum{TRUE}\hldef{)}
\hlkwd{head}\hldef{(y[,}\hlnum{1}\hlopt{:}\hlnum{9}\hldef{],}\hlnum{3}\hldef{)}
\end{alltt}
\begin{verbatim}
 Year Low High Hi_min Lo_max Avg_min Avg_max Mean Precip
 1900  -7   50     36     12      18    34.7 26.3   1.06
 1901  -7   48     37     20      17    31.8 24.4   1.45
 1902  -4   41     27     11      15    30.4 22.7   0.60
\end{verbatim}
\end{kframe}
\end{knitrout}

\bi

\item  We have to deal with the NA measurement for 1955. A simple approach is to replace the NA by the mean.

\item  What other approaches can you think of for dealing with this missing observation?

\item What are the strengths and weaknesses of these approaches?

\ei

\begin{knitrout}\small
\definecolor{shadecolor}{rgb}{0.969, 0.969, 0.969}\color{fgcolor}\begin{kframe}
\begin{alltt}
\hldef{low} \hlkwb{<-} \hldef{y}\hlopt{$}\hldef{Low}
\hldef{low[}\hlkwd{is.na}\hldef{(low)]} \hlkwb{<-} \hlkwd{mean}\hldef{(low,} \hlkwc{na.rm}\hldef{=}\hlnum{TRUE}\hldef{)}
\end{alltt}
\end{kframe}
\end{knitrout}

\end{frame}

\subsection{Smoothing the periodogram}

\begin{frame}[fragile]

\begin{knitrout}\small
\definecolor{shadecolor}{rgb}{0.969, 0.969, 0.969}\color{fgcolor}\begin{kframe}
\begin{alltt}
\hlkwd{spectrum}\hldef{(low,} \hlkwc{main}\hldef{=}\hlsng{"Unsmoothed periodogram"}\hldef{)}
\end{alltt}
\end{kframe}

{\centering \includegraphics[width=\maxwidth]{tmp/figure/periodogram-1} 

}


\end{knitrout}

\bi
\item  To smooth, we use the default periodogram smoother in {\Rlanguage}
\ei

\end{frame}


\begin{frame}[fragile]

\vspace{-3mm}

\begin{knitrout}\small
\definecolor{shadecolor}{rgb}{0.969, 0.969, 0.969}\color{fgcolor}\begin{kframe}
\begin{alltt}
\hlkwd{spectrum}\hldef{(low,} \hlkwc{spans}\hldef{=}\hlkwd{c}\hldef{(}\hlnum{3}\hldef{,}\hlnum{5}\hldef{,}\hlnum{3}\hldef{),} \hlkwc{main}\hldef{=}\hlsng{"Smoothed periodogram"}\hldef{,}
  \hlkwc{ylim}\hldef{=}\hlkwd{c}\hldef{(}\hlnum{15}\hldef{,}\hlnum{100}\hldef{))}
\end{alltt}
\end{kframe}
\end{knitrout}

\vspace{-4mm}

\begin{knitrout}\small
\definecolor{shadecolor}{rgb}{0.969, 0.969, 0.969}\color{fgcolor}

{\centering \includegraphics[width=10cm]{tmp/figure/smoothed_periodogram-1} 

}


\end{knitrout}
\bi
\item The bar is a 95\% pointwise confidence interval which we can slide to any frequency of interest.
\item The chi-squared property (\ref{eq:cs:log-chi-squared}) means this CI is the same width for each frequency, on the log scale. Note it is asymmetric.
\ei
\end{frame}


\begin{frame}[fragile]

\myquestion.  What is the default periodogram smoother in {\Rlanguage}?

\answer{\vspace{20mm}}{todo}

\myquestion.  How should we use it?

\answer{\vspace{20mm}}{todo}

\myquestion. Why is that default chosen?

\answer{\vspace{20mm}}{todo}

\end{frame}

\subsection{Tapering before calculating the  periodogram}

\begin{frame}[fragile]
\frametitle{More details on computing and smoothing the periodogram}
\bi
\item  To see what {\Rlanguage} actually does to compute and smooth the periodogram, type \code{?spectrum}. 

\item  This will lead you to type \code{?spec.pgram}.

\item  You will see that, by default, {\Rlanguage} removes a linear trend, fitted by least squares. This may often be a sensible thing to do. Why?

\item  You will see that {\Rlanguage} then multiplies the data by a quantity called a \myemph{taper}, computed by \code{spec.taper}. 

\item The taper smooths the ends of the time series and removes high-frequency artifacts arising from an abrupt start and end to the time series.

\item Formally, from the perspective of the Fourier transform, the time series takes the value zero outside the observed time points $1{\mycolon}N$. The sudden jump to and from zero at the start and end produces unwanted effects in the frequency domain.

\ei

\end{frame}

\begin{frame}[fragile]

The default taper in {\Rlanguage} smooths the first and last $p=0.1$ fraction of the time points, by modifying the detrended data $\data{y_{1:N}}$ to tapered version ${z_{1:N}}$ defined by
\begin{equation} \nonumber
{z_n} =
  \left\{
  \begin{array}{ll}
    \data{y_n} \big(1-\cos(\pi n/Np)\big)/2 & \mbox{ if $1\le n< Np$ } \\
    \data{y_n}  & \mbox{ if $Np \le n \le N(1-p)$ } \\
    \data{y_n} \big(1-\cos(\pi [N+1-n]/Np)\big)/2 & \mbox{ if $N(1-p)<n\le N$ }
  \end{array}\right.
\end{equation}

\vspace{-3mm}

\begin{knitrout}\small
\definecolor{shadecolor}{rgb}{0.969, 0.969, 0.969}\color{fgcolor}\begin{kframe}
\begin{alltt}
\hlkwd{plot}\hldef{(}\hlkwd{spec.taper}\hldef{(}\hlkwd{rep}\hldef{(}\hlnum{1}\hldef{,}\hlnum{100}\hldef{)),}\hlkwc{type}\hldef{=}\hlsng{"l"}\hldef{,}
  \hlkwc{main}\hldef{=}\hlsng{"Default taper in R, for a time series of length 100"}\hldef{)}
\hlkwd{abline}\hldef{(}\hlkwc{v}\hldef{=}\hlkwd{c}\hldef{(}\hlnum{10}\hldef{,}\hlnum{90}\hldef{),}\hlkwc{lty}\hldef{=}\hlsng{"dotted"}\hldef{,}\hlkwc{col}\hldef{=}\hlsng{"red"}\hldef{)}
\end{alltt}
\end{kframe}
\end{knitrout}

\vspace{-3mm}

\begin{knitrout}\small
\definecolor{shadecolor}{rgb}{0.969, 0.969, 0.969}\color{fgcolor}

{\centering \includegraphics[width=8cm]{tmp/figure/taper_plot-1} 

}


\end{knitrout}

\end{frame}


\subsection{Fitting an AR model to estimate the spectrum}

\begin{frame}[fragile]{Spectral density estimation by fitting a model}

Another standard way to estimate the spectrum is to fit an AR(p) model with $p$ selected by AIC.

\vspace{-2mm}

\begin{knitrout}\small
\definecolor{shadecolor}{rgb}{0.969, 0.969, 0.969}\color{fgcolor}\begin{kframe}
\begin{alltt}
\hlkwd{spectrum}\hldef{(low,}\hlkwc{method}\hldef{=}\hlsng{"ar"}\hldef{,}
  \hlkwc{main}\hldef{=}\hlsng{"Spectrum estimated via AR model picked by AIC"}\hldef{)}
\end{alltt}
\end{kframe}
\end{knitrout}

\vspace{-2mm}

\begin{knitrout}\small
\definecolor{shadecolor}{rgb}{0.969, 0.969, 0.969}\color{fgcolor}

{\centering \includegraphics[width=11cm]{tmp/figure/ar_periodogram-1} 

}


\end{knitrout}

\end{frame}

\begin{frame}[fragile]

\frametitle{Units of frequency and period}

\bi

\item When we call $\omega$ the frequency in cycles per unit time, we really mean \myemph{cycles per unit observation}.

\item  Suppose the time series consists of equally spaced observations, with $t_{n}-t_{n-1}=\Delta$ years. Then, the frequency is $\omega/\Delta$ \myemph{cycles per year}. 

\item  The \myemph{period} of an oscillation is the time for one cycle,
\begin{equation} \mbox{period} = \frac{1}{\mbox{frequency}}.\end{equation}

\item  When the observation intervals have a time unit (years, seconds, etc) we usually use that unit for the period, and its inverse for the frequency.

\ei

\end{frame}

\begin{frame}{Further reading} 

\bi

\item Sections~4.1 to 4.3 of \citet{shumway17} cover similar topics to this chapter.

\ei


\end{frame}


\newcommand\acknowledgments{
\begin{itemize}
\item   Compiled on {\today} using \Rlanguage version 4.4.2.
\item   \parbox[t]{0.75\textwidth}{Licensed under the \link{http://creativecommons.org/licenses/by-nc/4.0/}{Creative Commons Attribution-NonCommercial license}.
    Please share and remix non-commercially, mentioning its origin.}
    \parbox[c]{1.5cm}{\includegraphics[height=12pt]{../cc-by-nc}}
\item We acknowledge \link{https://ionides.github.io/531w24/acknowledge.html}{previous versions of this course}.
\end{itemize}
}

\mode<presentation>{
\begin{frame}[allowframebreaks=0.8]{References and Acknowledgements}
   
\bibliography{../bib531}

\vspace{3mm}

\acknowledgments

\end{frame}
}

\mode<article>{

{\bf \Large \noindent Acknowledgments}

\acknowledgments

  \bibliography{../bib531}

}

\end{document}


